\section{Reviewable Core Specification}\label{sec:appendix:specification}

This appendix makes the normative state machine, authority, event, recovery,
attribution, repair, and final-manuscript checks inspectable without relying
on source-file names. Any transition or permission not listed here is denied.

\subsection{Complete component-lifecycle transitions T1--T21}\label{sec:appendix:transitions}

The ten states are candidate, validated, shadow, probationary active, active,
warning, probation, quarantine, retired, and banned. Every rule commits in a
boundary transaction; T16 creates an immediate emergency boundary.

\small
\begin{longtable}{@{}p{0.045\textwidth}p{0.17\textwidth}p{0.31\textwidth}p{0.40\textwidth}@{}}
  \caption{Complete component-lifecycle transition table. Unlisted edges
  are illegal.}\label{tab:component-transitions}\\
  \toprule
  \textbf{Rule} & \textbf{Edge} & \textbf{Trigger} & \textbf{Guard and effect} \\
  \midrule
  \endfirsthead
  \toprule
  \textbf{Rule} & \textbf{Edge} & \textbf{Trigger} & \textbf{Guard and effect} \\
  \midrule
  \endhead
  T1 & candidate$\rightarrow$validated & candidate validation passes &
    Candidate cap is not exceeded and constitutional constraints are present.\\
  T2 & candidate$\rightarrow$retired & validation fails or candidate expires &
    Rejection is restricted to a candidate that never served.\\
  T3 & validated$\rightarrow$shadow & replay and anchor evaluation pass &
    A shadow slot is available.\\
  T4 & shadow$\rightarrow$validated & shadow sample is insufficient &
    Retry count remains within its limit.\\
  T5 & shadow$\rightarrow$retired & shadow quality is below threshold &
    Rejection is restricted to a candidate that never served.\\
  T6 & shadow$\rightarrow$probationary active & agreement threshold passes &
    The role is open; at most one adoption per role and boundary.\\
  T7 & probationary active$\rightarrow$active & clean probation record &
    Minimum probation service is met.\\
  T8 & probationary active$\rightarrow$quarantine & upheld impeachment or reliability collapse &
    Applied at an ordinary boundary.\\
  T9 & active$\rightarrow$warning & reliability warning or low-severity attack &
    The target continues serving.\\
  T10 & active$\rightarrow$probation & repeated warning or medium attack &
    The target continues serving.\\
  T11 & active$\rightarrow$quarantine & upheld impeachment or critical attack &
    A fallback is assigned.\\
  T12 & warning$\rightarrow$active & one clean epoch & Clear the warning.\\
  T13 & warning$\rightarrow$probation & warning recurs within memory &
    The target continues serving.\\
  T14 & probation$\rightarrow$active & required clean epochs complete &
    Minimum probation service is met.\\
  T15 & probation$\rightarrow$quarantine & degradation continues or impeachment is upheld &
    Applied at an ordinary boundary.\\
  T16 & serving state$\rightarrow$quarantine & critical fixed-kernel violation &
    From probationary active, active, warning, or probation; requires attached
    evidence, exactly one sanction, and a forced close/open boundary.\\
  T17 & quarantine$\rightarrow$retired & retirement adjudication confirmed &
    Replay confirmation is required and a retirement record is appended.\\
  T18 & quarantine$\rightarrow$probationary active & exoneration &
    Re-entry begins under probation.\\
  T19 & retired$\rightarrow$banned & contamination or critical finding &
    Banned is absorbing.\\
  T20 & active$\rightarrow$retired & adopted successor supersedes incumbent &
    Utility-policy role only; append a retirement record.\\
  T21 & active$\rightarrow$shadow & adopted successor supersedes incumbent &
    Paper writer/reviewer only; retain the old version as reinstatable standby.\\
  \bottomrule
\end{longtable}
\normalsize

T16 no longer changes a serving set inside an unchanged epoch. It force-closes
the current epoch and atomically commits the quarantine, old-epoch close, and
new-epoch open under a new fingerprint. The prior utility policy remains
frozen exactly; no successor is adopted in the emergency transaction. Samples
before and after T16 therefore have distinct institutional identities.

\paragraph{Epoch identity.}
The composite epoch fingerprint contains two separately reported
digests. The \emph{policy fingerprint} commits active component and prompt
versions, every registered prompt hash, the utility policy, registry
version, constitutional hash, transition thresholds, budgets, and other
resolved governance settings. The \emph{execution fingerprint} commits
the declared model backend and identifier, temperature, search and
evaluation settings, disabled tools, skill set, and explicit pins for model
revision, tool bundle, environment, and data snapshot. An absent external
revision is recorded as
\emph{unresolved}, not silently treated as fixed. Equality of policy
fingerprints therefore identifies equal declared governance policy;
equality of complete execution fingerprints identifies equal declared
execution dependencies. Neither proves equality of opaque provider
weights.

\paragraph{Boundary resolution order.}
The resolver freezes its reports, evaluations, registries, and settings,
then processes registered components in identifier order before unpaired
prompt candidates in identifier order. Component sanctions establish role
openings. Each object takes at most one lifecycle edge. T6 additionally
requires an opening and permits at most one adoption per role and boundary.
For a qualifying utility-policy successor, T20 retirement of its incumbent
and T6 adoption are placed in the same transaction; for a paper-role
successor, T21 moves the incumbent to shadow in that transaction. Events
are serialized in the fixed group order adoptions, sanctions, retirements,
and bans, but no intermediate group is visible without the transaction
commit. The post-state active map is then derived and each removed serving
role receives the latest eligible previous holder or a fixed baseline.
Kernel validation precedes the single atomic boundary commit.

\begin{table}[htbp]
  \centering
  \small
  \begin{tabularx}{\textwidth}{@{}p{0.34\textwidth}X@{}}
    \toprule
    \textbf{Whole-state invariant} & \textbf{Enforcement status} \\
    \midrule
    At most one primary serving holder per role &
      deterministic post-state projection and T6 per-role cap \\
    One serving set per composite fingerprint &
      boundary-only updates; T16 creates a fresh epoch \\
    Candidate authority cannot exceed incumbent authority &
      fixed capability-subset check before adoption \\
    Unadjudicated attacks cannot sanction or score &
      validated-attack requirement in kernel and scoring path \\
    Retired-dependent records cannot re-enter selection &
      repair closure plus eligibility check \\
    Uncommitted transactions do not replay &
      prepare/member/commit projection rule \\
    No post-T16 sample uses the old epoch identity &
      emergency close/quarantine/open transaction \\
    \bottomrule
  \end{tabularx}
  \caption{Composed-state invariants enforced by the implementation. They
  are regression tested but have not yet been model checked or proved.}
  \label{tab:whole-state-invariants}
\end{table}

\subsection{Authority matrix}\label{sec:appendix:authority}

The closed action vocabulary is read, append, write, invoke, and activate.
The resources are records, registry, audit log, frontier, active prompt text,
retired prompt text, checkpoint artefacts, and candidates. Unlisted pairs are
denied, and retired prompt text is unreadable by every role.

\begin{table}[htbp]
  \centering
  \scriptsize
  \setlength{\tabcolsep}{3.5pt}
  \begin{tabularx}{\textwidth}{@{}p{0.19\textwidth}*{8}{>{\centering\arraybackslash}X}@{}}
    \toprule
    \textbf{Role--tier row} & \textbf{Rec.} & \textbf{Reg.} &
    \textbf{Audit} & \textbf{Front.} & \textbf{Active} &
    \textbf{Retired} & \textbf{Artef.} & \textbf{Cand.} \\
    \midrule
    Evolvable--institutional & RAI & -- & -- & -- & R & -- & R & -- \\
    Evolvable--meta & RA & -- & -- & -- & R & -- & -- & -- \\
    Utility policy--institutional & RA & -- & -- & -- & -- & -- & -- & -- \\
    Governance actor--institutional & RAI & -- & RA & -- & R & -- & R & -- \\
    Transition engine--fixed & R & RW & RA & R & -- & -- & R & A \\
    Kernel/verifier--fixed & R & R & R & R & -- & -- & R & -- \\
    Audit log--fixed & -- & -- & A & -- & -- & -- & -- & -- \\
    \bottomrule
  \end{tabularx}
  \caption{Complete sparse authority matrix by resource. R, A, W, and I mean
    read, append, write, and invoke; A in the candidate column means activate.
    Evolvable rows cover every listed evolvable role at that tier; governance
    actors are auditor, evidence clerk, and governance judge. All omitted
    cells and all retired-prompt reads are denied. The implementation also
    emits the fully expanded role--tier--action--resource decision manifest.}
  \label{tab:authority-matrix}
\end{table}

A candidate's grants must be a subset of the incumbent grants it would
replace. Runtime typed builders ignore author or role fields in model text and
stamp the active component identifier and fixed role constant. This prevents
raw model output from self-asserting an identity at the application boundary.
It is not cryptographic or operating-system isolation: roles share a process,
there are no per-role keys or OS users, and checkpoint and tool access are not
separated by an IPC sandbox. The runtime, checkpoint permissions, and tool
gateway are consequently part of the trusted computing base.

\begin{table}[htbp]
  \centering
  \small
  \begin{tabularx}{\textwidth}{@{}p{0.28\textwidth}XX@{}}
    \toprule
    \textbf{Attacker capability} & \textbf{Presently enforced} &
    \textbf{Not guaranteed} \\
    \midrule
    Malicious text from a governed model &
      typed author stamping, API allow-list, legal transitions &
      semantic truth of the model's evidence \\
    Colluding or correlated model roles &
      distinct records and self-adjudication rejection &
      statistical independence or collusion resistance \\
    Buggy or compromised in-process role code &
      checks reached through the normal gateway &
      prevention of direct memory or file access \\
    Checkpoint writer or host administrator &
      no adversarial guarantee &
      registry, log, pin, or executable integrity \\
    Provider changes behind a stable model name &
      incompleteness recorded in execution identity &
      equality of opaque model weights \\
    External irreversible side effect &
      no fixed side-effect gate in this artifact &
      fail-closed device, API, or production-data safety \\
    \bottomrule
  \end{tabularx}
  \caption{Threat-capability matrix. The present claim is application-level
  reference-monitor and workflow integrity under a trusted host, not an
  operating-system or cryptographic security boundary.}
  \label{tab:threat-capability}
\end{table}

\subsection{Event envelope and recovery}\label{sec:appendix:events}

A new state-changing event uses schema version two and carries a monotonic
identifier, type, transaction identifier, payload, predecessor digest, and
timestamp. Its full SHA-256 digest commits the canonical tuple
\[
  (\textit{schema version},\textit{event id},\textit{event type},
  \textit{transaction id},\textit{payload},\textit{predecessor digest}).
\]
Thus type changes, identifier changes, transaction-boundary changes,
predecessor rewrites, payload edits, and reordering change the digest.
Timestamps remain metadata outside the digest. Legacy version-one
payload-only events remain readable, but new writers emit only version two.
This digest is not an author signature.

\begin{algorithm}[htbp]
\caption{Deterministic resume projection}
\begin{algorithmic}[1]
\Require Stored event sequence $E$ and, if it survives, local boundary pin $P$
\Ensure Reconstructed epoch and registries, or an integrity violation
\State Check strictly increasing identifiers
\State Recompute each version-aware event digest and predecessor link
\If{$P$ exists}
  \State Require $|E|\geq P.length$ and the hash at $P.length$ to equal $P.head$
\EndIf
\State Require matching transaction identifiers on prepare, members, and commit
\State Discard a trailing prepare sequence with no matching commit
\State Fold committed events from empty state to reconstruct epoch and registries
\State \Return reconstructed state
\end{algorithmic}
\end{algorithm}

Version two detects edits to replay-relevant fields, event reordering,
transaction-boundary substitution, interior deletion with surviving
successors, and a log shorter than a surviving local pin. Writers refuse to
extend malformed chains; replay ignores only an unterminated final
transaction as crash recovery. The design still cannot detect rollback of
both the log and its local pin to the same older prefix. It has no external
signed head, monotonic hardware counter, or independent witness.
Whole-checkpoint rollback is therefore outside the present threat model.

\subsection{Attribution, impact repair, and manuscript verification}\label{sec:appendix:algorithms}

\paragraph{Attribution.}
Raw attacks target artefacts. After adjudication, a fixed case-type table
selects implicated roles and resolves each against the epoch-frozen
incumbent. Each validated-attack record names one accountable component; a
multi-role case creates one record per role referencing the same judgment.
The seven general research attacks name the generator only when the attacked
node carries write-once producer component, prompt, and epoch provenance
matching the epoch-frozen generator. Legacy, missing, or mismatched provenance
remains targetless rather than guessing. Paper-reviewer over-acceptance binds
the reviewer and, when the draft also fails the fixed faithfulness gate, its
writer. Multiple causes are represented by multiple single-target records
sharing one judgment. The causal accuracy of these fixed mappings is
unevaluated; the one-target record shape is an audit constraint, not a causal
guarantee, and tool, data, or shared-prompt causes are not yet represented.

\paragraph{Impact repair.}
Direct impact consists of records stamped with a retired prompt hash.
Transitive impact is a breadth-first traversal of reverse source-reference
edges. Only proposal-to-proposal inspiration is declared non-material;
unknown edge types default to material. The traversal still computes the full
reachable closure after the default materiality depth of eight; the exemption
for a declared non-material edge stops at that depth, so deeper reachable
records are conservatively invalidated. With $V$ reachable records and $E$
reverse references, time and memory are $O(V+E)$, and a visited set terminates
cycles. Direct
generator/router/policy impact invalidates a node; reviewer/adversary/defender/
judge impact recomputes from surviving inputs and invalidates when those
inputs are insufficient. Semantic precision and over-removal remain
unevaluated.

\paragraph{Final manuscript verification.}
The verifier is not a general natural-language reasoner. The manuscript
forward-declares claim and numeric identifiers, formula, operand node/metric
locations, unit, and tolerance. The verifier checks executed-node and
artefact existence, recomputes a closed formula vocabulary, and compares the
reported value under absolute or relative tolerance. Strict scanning covers
the abstract, results, and conclusion; introduction, discussion, and
limitations produce warnings; related work, references, appendices, and
equations are excluded. It reports both reproducible-claim rate and numeric
coverage rate, plus identifier collisions, unknown formulas, missing
operands, units, and environment mismatch. This report was not produced from
an ARI execution checkpoint, so no valid execution-grounded claim manifest
exists for it and self-application has not been claimed. Substituting a
hand-made checkpoint would not test the proposed release path. Unstructured non-numeric truth,
the truthfulness of execution records, and semantic paraphrase are outside
scope. ``Fixed'' means model-free determinism over this structured input.

\subsection{Boundary of determinism and test evidence}\label{sec:appendix:determinism}

We claim only \emph{kernel determinism}: saved advisory reports, candidate
evaluations, registries, and configuration produce the same transition and
repair; and \emph{log-replay determinism}: the same committed event sequence
produces the same projection. We do not claim that a model regenerates the
same defence, adjudication proposal, or review from the same prompt.

At this revision, the complete suite collected 4,805 tests: 4,787 passed,
16 were skipped, and two were expected failures. All 1,124 tests directly
covering RQGM and the final claim gate passed. This revision adds negative
cases for semantic event-field edits, predecessor rewrites, ordering,
transaction mismatch, emergency-boundary identity, and authority-manifest
completeness. The non-overlapping grouped breakdown is 64 claim gate; 285
kernel/transition/epoch/state/founding/mode; 126 adversarial/governance;
66 repair/clean room; 85 utility evolution; 157 evaluation; 93 prompt
evolution; 94 paper roles; and 154 other focused tests. Counts are not
performance evidence.
Code coverage, mutation score, multi-model detection rates, cost, and
research-quality comparisons have not been measured.
